\section{Related Works}



\textbf{Fine-tuning Large Language Models:} Large language models (LLMs), including GPT-4, PaLM, and LLaMa, have demonstrated exceptional proficiency in a variety of natural language assignments. Recent progress in instruction tuning has enabled LLMs to adeptly comprehend user directives and generate responses that resemble human-like communication. (\cite{hu2024finetuning}) proposed sequential instructions mechanism for fine-tuning LLMs, which is an effective strategy to automatically augment instruction tuning data and equip LLMs with the ability to execute multiple sequential instructions. (\cite{doan2023efficient}) leverage large-scale instruction-following datasets from open-source projects like Alpaca, GPT4All, and Chat-Doctor, which cover general domain and specific medical domain.

\textbf{Retrieval-Augmented Generation:} In recent years, Large Language Models (LLMs) have emerged as the cornerstone of natural language processing, showcasing unparalleled proficiency in tasks ranging from text generation to translation and summarization. However, despite their remarkable capabilities, LLMs come with inherent limitations, particularly in their ability to access and integrate external knowledge effectively. This drawback often manifests in the form of generating contextually inaccurate or incomplete responses, especially when confronted with complex queries or specialized domains. (\cite{lewis2020retrieval}) propose a general-purpose finetuning recipe for retrieval augmented generation (RAG). Building upon the foundation laid by LLMs, RAG revolutionizes language generation by enabling dynamic access to a vast array of external knowledge sources, such as databases, documents, or online repositories.

\textbf{Multimodal Retrieval-Augmented Generation:} In the ever-evolving landscape of artificial intelligence, the integration of multimodal capabilities has emerged as a pivotal frontier, enabling machines to comprehend and generate content across multiple modalities, including text, images, and videos. Amidst this paradigm shift, Multimodal Retrieval Augmented Generation (MRAG) emerges as a groundbreaking approach that marries the power of multimodal understanding with the sophistication of retrieval-augmented generation techniques.
Recent MRAG (\cite{chen2022murag}, \cite{yasunaga2022retrieval}) that seeks to overcome these limitations by amalgamating multimodal understanding with dynamic information retrieval mechanisms. By harnessing the complementary strengths of multimodal AI models and retrieval-augmented generation techniques, MRAG enables machines to access a vast reservoir of external knowledge sources spanning textual, visual, and auditory domains.

\textbf{Single and Multi Representations in Dense Passage Retrieval:} ColBERT (Contextualized Late Interaction over BERT) (\cite{khattab2020colbert}) is a dense retrieval approach that utilizes multiple representations for passage indexing and retrieval. Unlike single representation methods, such as ANCE, which represent entire passages with a single embedding (usually BERT's [CLS] token), ColBERT assigns an embedding to each token in a passage. This approach allows ColBERT to capture fine-grained information and context within passages, potentially leading to improved retrieval effectiveness. By representing passages at a more granular level, ColBERT can better match queries with relevant passages, especially for complex information needs and definitional queries.